\subsection{正则序列与纯量扩张}

命题 5.--- 设 \(\mathfrak{p}: A \to B\) 为 Noether 局部环的局部同态，N 为有限生成 B-模，\(\mathbf{x} = (x_1, \ldots, x_r)\) 为 \(\mathfrak{m}_B\) 的元素的一个序列，\(u: A[T_1, \ldots, T_r] \longrightarrow B\) 为唯一的 A-代数同态，使 \(u(T_i) = x_i\)（对 \(i = 1, \ldots, r\)）。下列条件等价：

\begin{enumerate}
\def\labelenumi{(\roman{enumi})}
\item
  同态 u 使 N 成为平坦的 A{[}T \(_{1}\) ,\ldots{} ,T \(_{r}\) {]}-模；
\item
  A-模 N 是平坦的，且对每个 A-模 M，序列 x 是 M ⊗A N-正则的；
\item
  A-模 N 是平坦的且序列 x 是 \(\kappa_{A} \otimes_{A} N\)-正则的；
\item
  A-模 N/(x₁N + \ldots{} + xᵣN) 是平坦的且序列 x 是 N-正则的。
\end{enumerate}

(i)⇒(ii)：令 \(\mathbf{T} = (\mathrm{T}_{1}, \ldots, \mathrm{T}_{r})\)。设 A{[}T{]}-模 N 是平坦的。由于 A{[}T{]} 在 A 上平坦，A-模 N 是平坦的 (I, § 2, 第 7 节, 命题 8 的推论 3)。设 M 为 A-模。序列 T 显然是 M{[}T{]}-正则的，从而是 M{[}T{]} ⊗\_\{A{[}T{]}\} N-正则的，因为 N 在 A{[}T{]} 上平坦。而 A{[}T{]}-模 M{[}T{]} ⊗\_\{A{[}T{]}\} N 等同于 M ⊗\_\{A\} N，比为 T\_i 的同态对应于自同态 1\_M ⊗ (x\_i)\_N。因此条件 (ii) 满足。

\begin{enumerate}
\def\labelenumi{(\roman{enumi})}
\setcounter{enumi}{1}
\tightlist
\item
  \(\Rightarrow\) (iii)：这是显然的。
\end{enumerate}

(iii)⇒(iv)：这由 § 1, 第 6 节的命题 10 得出。

(iv)⇒(i)：用 t 表示 A{[}T{]} 中由 T 生成的理想。(A{[}T{]}/t)-模 N/tN 按假设是平坦的，且 N 对 t 是理想分离的 (III, § 5, 第 4 节, 命题 2)。为证明 N 在 A{[}T{]} 上平坦，鉴于同上引文第 2 节的定理 1，只需证明 A-模 \(\operatorname{Tor}_{1}^{A{[}T{]}}(A{[}T{]}/t,N)\) 为零。但由于序列 T 是 A{[}T{]}-正则的，该模同构于 \(\mathrm{H}_{1}(T,N)\) (A, X, 第 159 页, 注 3)，而它为零，因为序列 T 是 N-正则的（同上引文, 第 157 页, 命题 5）。

推论. 设 \(k\) 为域，\(\Lambda\) 为一个 \(k\) 上的 Noether 局部代数，\(\mathbf{x} = (x_1, \ldots, x_r)\) 为 \(\mathfrak{m}_{\mathrm{A}}\) 的元素的一个序列，M 为 \(\Lambda\) 上的有限型模。用 \(\widehat{\mathrm{A}}\) 和 \(\widehat{\mathrm{M}}\) 表示 A 和 M 关于其 ( \(x_1\mathrm{A} + \ldots + x_r\mathrm{A}\) -进拓扑的完备化；用 \(u: k[\mathrm{T}_1, \ldots, \mathrm{T}_r] \longrightarrow \mathrm{A}\) 表示唯一的 \(k\) -代数同态，使 \(u(\mathrm{T}_i) = x_i\)（对 \(i = 1, \ldots, r\)），并用 \(\hat{u}: k[[\mathrm{T}_1, \ldots, \mathrm{T}_r]] \longrightarrow \widehat{\mathrm{A}}\) 表示扩张它的唯一连续同态。下列条件等价：

\begin{enumerate}
\def\labelenumi{(\roman{enumi})}
\item
  序列 x 是 M-正则的；
\item
  同态 u 使 M 成为平坦的 \(k[T_{1},\ldots,T_{r}]\) -模；
\item
  同态 \(\hat{u}\) 使 \(\hat{M}\) 成为平坦的 \(k[[T_{1},\ldots,T_{r}]]\) -模。
\end{enumerate}

(i) 与 (ii) 的等价性由命题 5 的条件 (i) 与 (iv) 的等价性得出；(ii) 与 (iii) 的等价性由 III, § 5, 第 4 节, 命题 4 得出。

这些结果使我们能在两个重要情形刻画 Macaulay 模。用 \(\Lambda\) 表示一个 Noether 局部环，M 为一个有限生成 A-模。说 \(\Lambda\) -模 M 是 Macaulay 的等价于说 \(\widehat{A}\) -模 \(\widehat{M}\) 是 Macaulay 的 ( \(\S\) 2, 第 7 节, 命题 8 的推论 4)。以下我们总假设 Noether 局部环 A 是完备的。

\begin{enumerate}
\def\labelenumi{\arabic{enumi})}
\tightlist
\item
  先设 A 具有一个子域；于是它具有一个代表域 \(k\) (IX, § 3, 第 3 节, 定理 1)。设 \((x_{1},\ldots ,x_{r})\) 为对 M 的一个极大割序列；用 \(u:k[[\mathrm{T}_{1},\dots ,\mathrm{T}_{r}]]\longrightarrow \Lambda\) 表示唯一的连续同态，使 \(u(\mathrm{T}_i) = x_i\)（对 \(i = 1,\dots ,r\)）。由 IX, § 2, 第 5 节的引理 4 b) 及 VIII, § 3, 第 2 节的注 1，A/Ann(M) 是有限型 \(k[[\mathrm{T}_{1},\dots ,\mathrm{T}_{r}]]\) -模，从而 M 是有限型 \(k[[\mathrm{T}_{1},\dots ,\mathrm{T}_{r}]]\) -模。这样，下列条件等价：
  \begin{enumerate}
  \def\labelenumii{(\roman{enumii})}
  \item
    \(k[[\mathrm{T}_1,\dots ,\mathrm{T}_r]]\) -模 M 是自由的；
\item
    A-模 M 是 Macaulay 的。
\end{enumerate}

事实上，说 M 是 Macaulay A-模等价于说序列 \((x_{1},\ldots,x_{r})\) 是 M-正则的 (§ 2, 第 3 节, 定理 1)。由上面的推论，后一条件意味着模 M 在 \(k[[T_{1},\ldots,T_{r}]]\) 上平坦，等价地，由于它有限生成，它是自由的。

\item
  设 A 的剩余域 \(\kappa_{\Lambda}\) 的特征为 \(p > 0\)，且有 \(\dim(M/pM) < \dim(M)\) 。设 \((x_1, \ldots, x_r)\) 为对 M/pM 的一个极大割序列，于是 \((p1_{\mathrm{A}}, x_1, \ldots, x_r)\) 是对 M 的极大割序列。设 C 为一个 \(p\) -环，长度 \(+\infty\)，其剩余域为 \(\kappa_{\Lambda}\) (IX, § 2, 第 3 节, 命题 5)。存在一个从 C 到 A 中的同态 \(u_0\)，它在剩余域上诱导恒同；设 \(u: \mathrm{C}[[\mathrm{T}_1, \ldots, \mathrm{T}_r]] \longrightarrow \mathrm{A}\) 为扩张 \(u_0\) 且把 \(\mathrm{T}_i\) 映到 \(x_i\)（对每个 \(i\)）的唯一同态。与前面一样，由同上引文第 5 节的引理 4 得出，\(u\) 使 M 成为有限型 \(\mathrm{C}[[\mathrm{T}_1, \ldots, \mathrm{T}_r]]\) -模。下列条件等价：
  \begin{enumerate}
  \def\labelenumii{(\roman{enumii})}
\item
    C{[}{[}T1, \ldots, Tr{]}{]}-模 M 是自由的；
\item
    A-模 M 是 Macaulay 的。
\end{enumerate}

事实上，条件 (ii) 等价于说序列 \((x_{1},\ldots,x_{r})\) 是 M-正则的，且比为 p 的同态在 \(\mathrm{M}/(x_{1}\mathrm{M}+\ldots+x_{r}\mathrm{M})\) 中是单射的 (§ 2, 第 3 节, 定理 1)。但后一条件意味着 C-模 \(\mathrm{M}/(x_{1}\mathrm{M}+\ldots+x_{r}\mathrm{M})\) 是无挠的，从而是平坦的 (Λ, X, 第 9 页, 例 7)。于是，鉴于命题 5 的 (iv)⇔(i)，条件 (ii) 等价于模 M 在 \(C[T_{1},\ldots,T_{r}]\) 上平坦，或等价地 (III, § 5, 第 4 节, 命题 4)，\(C[[T_{1},\ldots,T_{r}]]\) -模 M 是平坦的，即由于它有限生成而是自由的。

\end{enumerate}

